\lysh{14652}{4}

\begin{alist}
\item Είναι:
\[x^2 > x \Leftrightarrow x^2 - x > 0 \quad (1)\]

Το τριώνυμο $x^2 - x = x(x - 1)$ έχει ρίζες τις: $x = 0$ και $x = 1$.
Το πρόσημο του τριωνύμου φαίνεται στον παρακάτω πίνακα.

\begin{center}
\begin{tabular}{|c|ccccccc|}

$x$ & $-\infty$ & & $0$ & & $1$ & & $+\infty$ \\ 
$x^2 - x$ & & $+$ & $0$ & $-$ & $0$ & $+$ & \\ 
\end{tabular}
\end{center}

Επομένως η $(1)$ αληθεύει για $x \in (-\infty, 0) \cup (1, +\infty)$.

\item 
\begin{rlist}
\item Από τα δεδομένα και το $\alpha)$ ερώτημα έχουμε: $0 < 1 < \alpha < \alpha^2$ και
$1 < \alpha$, οπότε $\dfrac{1}{\alpha} < \dfrac{\alpha}{\alpha}$, δηλαδή
$\dfrac{1}{\alpha} < 1$.
Επίσης: $\alpha < \alpha^2 \Leftrightarrow \dfrac{\alpha}{\alpha^2} < \dfrac{\alpha^2}{\alpha^2} \Leftrightarrow \dfrac{1}{\alpha} < 1$.
Οπότε τελικά: $0 < \dfrac{1}{\alpha} < 1 < \alpha < \alpha^2$.

\item Έχουμε διαδοχικά:
\begin{align*}
\alpha < \alpha^2 \ &\Leftrightarrow \\
\dfrac{\alpha}{2} < \dfrac{\alpha^2}{2} \ &\Leftrightarrow \\
\dfrac{\alpha}{2} + \dfrac{\alpha}{2} < \dfrac{\alpha^2}{2} + \dfrac{\alpha}{2} \ &\Leftrightarrow \\
\alpha &< \dfrac{\alpha^2 + \alpha}{2}.
\end{align*}

Επίσης:
\begin{align*}
\alpha < \alpha^2 \ &\Leftrightarrow \\
\dfrac{\alpha}{2} < \dfrac{\alpha^2}{2} \ &\Leftrightarrow \\
\dfrac{\alpha^2}{2} + \dfrac{\alpha}{2} < \dfrac{\alpha^2}{2} + \dfrac{\alpha^2}{2} \ &\Leftrightarrow \\
\dfrac{\alpha^2 + \alpha}{2} &< \alpha^2.
\end{align*}

Οπότε τελικά: $\alpha < \dfrac{\alpha^2 + \alpha}{2} < \alpha^2$.
\end{rlist}
\end{alist}
